\seccionreporte{Implementación del microservicio}{sec:implementacion}

\propositoseccion{%
  Resume la estructura del repositorio y los módulos críticos. No pegue todo el código: cite rutas y fragmentos representativos.%
}

\subsection{Repositorio}

URL: \ReporteRepositorio{} (\url{https://github.com/AleoshGG/vivia-ai}).

El servicio sigue una arquitectura MVC + SOLID, con el microservicio de detección
de anomalías organizado por capas dentro de \texttt{src/anomaly\_detector\_api/}.

\begin{tabular}{@{}lp{9.5cm}@{}}
  \toprule
  Ruta & Contenido \\
  \midrule
  \texttt{src/anomaly\_detector\_api/} & Código del microservicio (controllers, usecases, services, persistence, messaging, models) \\
  \texttt{shared/} & Componentes transversales: autenticación, carga de modelos, salud, cola \\
  \texttt{config/} & Configuración tipada (\texttt{settings.py}) y \textit{logging} \\
  \texttt{models\_registry/anomaly/} & Artefactos serializados del modelo (modelo, \textit{scaler}, columnas, metadatos) \\
  \texttt{notebooks/} & Entrenamiento (\texttt{04\_model\_training.ipynb}), EDA y generación del dataset \\
  \texttt{datasets/} & Datos de entrenamiento (dataset sintético) \\
  \texttt{scripts/} & Registro del modelo en MLflow (\texttt{register\_anomaly\_model.py}) \\
  \texttt{alembic/} & Migraciones de base de datos \\
  \texttt{tests/} & Pruebas automatizadas (unitarias e integración) \\
  \texttt{infrastructure/docker/} & \texttt{Dockerfile} por servicio (API, MLflow, \textit{batch}) \\
  \texttt{compose.yml} & Orquestación de servicios \\
  \texttt{README.md} & Instalación y ejecución \\
  \bottomrule
\end{tabular}

\subsection{Arranque del servicio}

En producción, el servicio se despliega de forma automatizada mediante un flujo
de GitHub Actions (\texttt{deploy.yml}) que sincroniza el código al servidor y
levanta los contenedores. El arranque del microservicio de IA es
\textbf{escalonado}: primero se levanta MLflow, después se registra el modelo de
forma idempotente y, por último, se inicia el resto de los servicios. El
siguiente fragmento corresponde al bloque de despliegue de \texttt{vivia-ai}:

\begin{lstlisting}[style=bashstyle]
# ── 3. vivia-ai (Python AI) ──────────────────────────────────────
# Credencial del service account de GCS (artefactos de MLflow).
mkdir -p ~/secrets
echo '${{ secrets.GCS_SA_JSON }}' > ~/secrets/gcs-vivia.json
chmod 600 ~/secrets/gcs-vivia.json

cat > ~/vps/vivia-ai/.env << EOF
RABBITMQ_HOST=vivia-rabbitmq
RABBITMQ_PORT=5672
RABBITMQ_USER=${{ secrets.RABBITMQ_USER }}
RABBITMQ_PASSWORD=${{ secrets.RABBITMQ_PASSWORD }}
RABBITMQ_VHOST=${{ secrets.RABBITMQ_VHOST }}
INTERNAL_API_KEY=${{ secrets.INTERNAL_API_KEY }}
EXTERNAL_SERVICE_URL=http://vivia-app:8080
MLFLOW_ARTIFACT_ROOT=${{ secrets.MLFLOW_ARTIFACT_ROOT }}
MODELS_REGISTRY_HOST=${{ secrets.MODELS_REGISTRY_HOST }}
MLFLOW_TRACKING_URI=http://mlflow:5000
POSTGRES_USER=${{ secrets.AI_DB_USER }}
POSTGRES_PASSWORD=${{ secrets.AI_DB_PASSWORD }}
POSTGRES_DB=${{ secrets.AI_DB_NAME }}
DATABASE_URL=postgresql+asyncpg://${{ secrets.AI_DB_USER }}:${{ secrets.AI_DB_PASSWORD }}@${{ secrets.AI_DB_HOST }}:5432/${{ secrets.AI_DB_NAME }}?ssl=require
MLFLOW_BACKEND_STORE_URI=postgresql://${{ secrets.AI_DB_USER }}:${{ secrets.AI_DB_PASSWORD }}@${{ secrets.AI_DB_HOST }}:5432/${{ secrets.AI_DB_NAME }}?sslmode=require
GCS_CREDENTIALS_HOST=$HOME/secrets/gcs-vivia.json
EOF

cd ~/vps/vivia-ai
docker compose down --remove-orphans
docker compose build --no-cache

# Arranque escalonado: MLflow primero (la BD ya es RDS, externa),
# registramos el modelo (idempotente) y por ultimo levantamos todo.
docker compose up -d mlflow
docker compose run --rm anomaly-api python -m scripts.register_anomaly_model
docker compose up -d
docker image prune -f
\end{lstlisting}

Para ejecución local de desarrollo, el servicio se levanta con \textit{uvicorn}
tras instalar las dependencias con Poetry; la aplicación FastAPI se define en
\texttt{src/anomaly\_detector\_api/main.py}.

\subsection{Inicialización (\textit{lifespan})}

Al arrancar, la aplicación FastAPI ejecuta un \textit{lifespan} que prepara las
dependencias del servicio una sola vez: construye el \textit{engine} de base de
datos y el repositorio de inferencias, carga el modelo desde el Model Registry de
MLflow y lanza el consumidor de la cola de RabbitMQ en un hilo en segundo plano.

\begin{lstlisting}[style=pythonstyle]
# src/anomaly_detector_api/main.py (extracto)
@asynccontextmanager
async def lifespan(app: FastAPI):
    # Base de datos: engine + repositorio de inferencias
    engine          = build_engine()
    session_factory = build_session_factory(engine)
    await init_models(engine)
    app.state.inference_repository = InferenceRepository(session_factory)

    # Modelo desde el Model Registry de MLflow (una sola carga)
    model = AnomalyModel(settings.anomaly_model_name, settings.anomaly_model_stage)
    model.load()
    app.state.anomaly_model = model

    # Consumidor de la cola de RabbitMQ en hilo demonio
    consumer = AnomalyQueueConsumer(model=model)
    thread   = threading.Thread(target=consumer.run, daemon=True,
                                name="anomaly-queue-consumer")
    thread.start()
    yield
    consumer.stop()
    await engine.dispose()
\end{lstlisting}

\subsection{Fragmento representativo}

El núcleo de la inferencia reside en el caso de uso
\texttt{AnalyzePropertyUseCase} (\texttt{src/anomaly\_detector\_api/usecases/}
\texttt{analyze\_property.py}): construye las \textit{features}, ejecuta la
predicción del modelo en un \textit{thread pool} (para no bloquear el
\textit{event loop}), persiste la inferencia y notifica el veredicto al servicio
externo.

\begin{lstlisting}[style=pythonstyle]
# src/anomaly_detector_api/usecases/analyze_property.py (extracto)
async def execute(self, request: PropertyRequest) -> dict:
    loop     = asyncio.get_running_loop()
    features = build_features(request.draft, self._model.feature_cols)

    # predict() es sincrono (numpy/sklearn) -> thread pool
    is_anomaly, score = await loop.run_in_executor(
        None, self._model.predict, features
    )

    approved = not is_anomaly
    reason = (
        "Propiedad aprobada tras analisis de anomalias."
        if approved
        else f"Propiedad rechazada: anomalia detectada (score={score:.4f})."
    )

    # La inferencia se persiste ANTES del webhook: el registro no debe
    # depender del exito de la notificacion al servicio externo.
    await self._persist(request, is_anomaly, score, approved, reason, features)

    payload = AnalysisPayload(draft_id=request.draft.id,
                              approved=approved, reason=reason)
    status_code = await self._post_result(payload)   # webhook con reintentos
    return self._response(payload, status_code)
\end{lstlisting}

\subsection{Manejo de errores}

\begin{itemize}
  \item \texttt{403}: cabecera \texttt{X-Internal-API-Key} ausente o inválida
  (acceso restringido a clientes internos).
  \item \texttt{422}: entrada inválida (campos faltantes o tipos incorrectos),
  con detalle estructurado en JSON generado por Pydantic.
  \item \texttt{404}: consulta de una inferencia inexistente por su identificador.
  \item \texttt{500}: fallo interno; se registra en \textit{logs} sin filtrar
  datos sensibles.
  \item \textbf{Tolerancia a fallos del webhook}: la notificación al servicio
  externo reintenta con \textit{backoff} (\texttt{[2, 5, 10]} s) y, ante un fallo
  persistente, envía un resultado de respaldo (\texttt{approved=false}) para que
  la publicación se revise manualmente; un fallo de persistencia, por su parte,
  no aborta el análisis.
\end{itemize}

\subsection{Dependencias principales}

\begin{tabular}{@{}ll p{6.0cm}@{}}
  \toprule
  Librería & Versión & Rol \\
  \midrule
  \texttt{fastapi}        & \textasciicircum0.109.0 & Framework web ASGI \\
  \texttt{uvicorn}        & \textasciicircum0.27.0  & Servidor ASGI \\
  \texttt{pydantic}       & \textasciicircum2.5.3   & Validación de esquemas \\
  \texttt{pydantic-settings} & \textasciicircum2.1.0 & Configuración tipada \\
  \texttt{mlflow}         & \textasciicircum2.10.0  & Model Registry y empaquetado \texttt{pyfunc} \\
  \texttt{scikit-learn}   & (vía MLflow)            & \textit{Isolation Forest} y \texttt{StandardScaler} \\
  \texttt{google-cloud-storage} & \textasciicircum2.14.0 & Artefactos de modelos en GCS \\
  \texttt{pika}           & \textasciicircum1.3.2   & Cliente AMQP de RabbitMQ \\
  \bottomrule
\end{tabular}

\noindent Dependencias de desarrollo: \texttt{pytest}, \texttt{pytest-asyncio},
\texttt{pytest-cov}, \texttt{black}, \texttt{ruff} y \texttt{jupyter}. La gestión
de dependencias se realiza con Poetry (\texttt{pyproject.toml}).
